% =====================================================================
% Metadata-RAG Full-Track Report
% A Paper-Derived Reimplementation and Its Empirical Behaviour
% =====================================================================
\documentclass[11pt,a4paper]{article}

\usepackage[utf8]{inputenc}
\usepackage[T1]{fontenc}
\usepackage{lmodern}
\usepackage[margin=1in]{geometry}
\usepackage{amsmath,amssymb,amsfonts}
\usepackage{booktabs}
\usepackage{multirow}
\usepackage{longtable}
\usepackage{array}
\usepackage{tabularx}
\usepackage{siunitx}
\usepackage{graphicx}
\usepackage{xcolor}
\usepackage{hyperref}
\usepackage{url}
\usepackage{enumitem}
\usepackage{listings}
\usepackage{fancyhdr}
\usepackage{titlesec}
\usepackage{caption}
\usepackage{subcaption}
\usepackage{microtype}
\usepackage{parskip}
\usepackage{algorithm}
\usepackage{algpseudocode}

\hypersetup{
  colorlinks=true,
  linkcolor=blue!50!black,
  citecolor=blue!50!black,
  urlcolor=blue!50!black,
  pdftitle={Metadata-RAG Full-Track Report},
  pdfauthor={RAG Architecture Readiness Program},
}

\sisetup{
  round-mode=places,
  round-precision=4,
  detect-weight=true,
  detect-family=true,
}

\definecolor{codebg}{RGB}{248,248,248}
\definecolor{codekw}{RGB}{0,90,160}
\definecolor{codestr}{RGB}{160,50,20}
\definecolor{codecom}{RGB}{110,110,110}

\lstdefinelanguage{yaml}{
  keywords={true,false,null},
  keywordstyle=\color{codekw}\bfseries,
  commentstyle=\color{codecom}\itshape,
  stringstyle=\color{codestr},
  morecomment=[l]{\#},
}
\lstset{
  basicstyle=\ttfamily\small,
  backgroundcolor=\color{codebg},
  frame=single,
  rulecolor=\color{black!20},
  breaklines=true,
  showstringspaces=false,
  captionpos=b,
  columns=fullflexible,
  keepspaces=true,
}

\titleformat{\section}{\large\bfseries}{\thesection}{1em}{}
\titleformat{\subsection}{\normalsize\bfseries}{\thesubsection}{1em}{}

\pagestyle{fancy}
\fancyhf{}
\rhead{\small Metadata-RAG Full-Track Report}
\lhead{\small Paper-Derived Reimplementation}
\cfoot{\thepage}

\title{
  \vspace{-1.5cm}
  \rule{\linewidth}{0.4pt}\\[0.3cm]
  {\Large\bfseries Metadata-RAG:}\\[0.15cm]
  {\large A Paper-Derived Reimplementation of Metadata-Driven\\
   Multistage Retrieval, and Its Empirical Behaviour on Four\\
   Heterogeneous Long-Context Corpora}\\[0.2cm]
  \rule{\linewidth}{0.4pt}\\[0.4cm]
  \normalsize RAG Architecture Readiness Program\\
  \small Technical Report --- Metadata-RAG Track Closure
}
\author{}
\date{\today}

\begin{document}

\maketitle
\thispagestyle{empty}

% =====================================================================
\begin{abstract}
\noindent
We describe, freeze, implement, and empirically characterise
\emph{Metadata-RAG}, a paper-derived reimplementation of the
metadata-driven multistage retrieval architecture described in
arXiv:2510.24402. Unlike the program's other baselines, Metadata-RAG
has no foreign source: the paper is a preprint under major revision
that defines an algorithmic family but pins neither its metadata
schema, its query-metadata inference rules, nor its fusion
coefficients. We therefore designed, froze, and executed a full
project-owned specification as a strict comparator in the program's
controlled architectural track. This report documents every frozen
design decision, every controlled adaptation, the complete
implementation (eleven production modules, six operational scripts,
nineteen test modules, \num{342} passing tests), and an end-to-end
execution on four heterogeneous canonical corpora (CLERC, CUAD,
GreekBarBench, LaRA). Under a shared retrieval-only protocol
identical to the program's other tracks, Metadata-RAG achieves a
macro-averaged document \texttt{Recall@10} of \num{0.3651}, against
LongRAG's \num{0.3144} and My\_RAG's \num{0.4049}. The most
scientifically interesting finding is not the aggregate number but
its decomposition: under the frozen \texttt{min\_match\_ratio} of
\num{0.5}, the metadata-filtering stage failed to activate on three
of the four corpora (\num{99.9}\% fallback on CLERC, \num{97.6}\% on
CUAD, \num{92.3}\% on GreekBarBench, and \num{33.7}\% on LaRA). We
document the mechanism, the reasons, and the consequences with full
honesty, and we make no post-hoc adjustment of the frozen
specification. The report concludes with the architectural
observations that are directly usable in the program's comparative
publication.
\end{abstract}

\vfill
\noindent\textbf{Keywords.} metadata-driven retrieval; paper-derived
reimplementation; clean-room specification; frozen thresholds;
deterministic windowing; dense-only retrieval; hybrid retrieval;
retrieval-only evaluation; reproducibility; architectural
comparators.

\newpage
\tableofcontents
\newpage

% =====================================================================
\section{Introduction}\label{sec:intro}

\subsection{Position of this work in the program}
The RAG Architecture Readiness Program is a systematic,
retrieval-only comparative study of multiple retrieval architectures
under a shared controlled evaluation protocol. Its central research
question is:

\begin{quote}
\emph{Under identical embedder, query set, ground truth, and
evaluator, which retrieval architecture returns the most relevant
canonical documents from long-context corpora?}
\end{quote}

Metadata-RAG participates in the program as a
\emph{paper-derived comparator}. It is neither a native-fidelity
reproduction of an author-provided artefact (no such artefact exists
for this paper) nor a controlled adaptation of an executed baseline
whose architecture can be examined in source. Instead, its design
has been reimplemented from the paper's algorithmic description
under a frozen project-owned specification. Its fidelity class is
therefore \textsc{PAPER\_DERIVED\_COMPARATOR}, and its
controlled-track eligibility is what this report establishes.

\subsection{The question specific to this track}
The specific question answered by this report is threefold:

\begin{enumerate}[leftmargin=1.5em,itemsep=0.15em]
  \item Can a metadata-driven multistage retrieval architecture be
    defined, frozen, and executed end-to-end under the same
    controlled protocol as the program's other tracks, without
    permitting hidden tuning?
  \item What is that architecture's measured document-level
    retrieval quality on the program's canonical four-dataset scope?
  \item What does the decomposition of that measurement reveal about
    the architecture's behaviour and, in particular, about the
    interaction between its frozen design choices and the properties
    of the corpora it is asked to retrieve from?
\end{enumerate}

\subsection{Contributions}
\begin{enumerate}[leftmargin=1.5em,itemsep=0.15em]
  \item A fully specified, frozen, and validated clean-room
    specification for a paper-derived reimplementation, including
    four explicit amendments that resolve every point at which the
    paper's design was under-determined.
  \item An implemented and unit-tested architecture: eleven
    production modules, six operational scripts, nineteen test
    modules, \num{342} passing tests, all under strict type checking
    and lint discipline.
  \item A deterministic oversize-document windowing policy, needed
    because three of the four datasets contain documents that exceed
    the BGE-M3 input budget by one to two orders of magnitude.
  \item A complete end-to-end execution of the architecture on
    CLERC, CUAD, GreekBarBench, and LaRA, with per-query and
    macro-aggregated metrics.
  \item A \emph{measured, honest account} of the architecture's
    metadata-stage activation rate on each corpus, and the
    architectural interpretation of that rate.
  \item A three-track comparison against the program's LongRAG and
    My\_RAG tracks over the identical four datasets.
\end{enumerate}

\subsection{What this report is not}
It is not a claim of native-fidelity reproduction. The paper is a
preprint under major revision; no foreign source exists; no
paper-level effectiveness claim is made against this baseline. It is
also not a claim of superiority. We report the measured behaviour of
a frozen design under a shared protocol; the decomposition of that
behaviour is the substantive contribution.

% =====================================================================
\section{Paper Identity and Reimplementation Class}\label{sec:identity}

\subsection{The paper}
Metadata-RAG is derived from
\emph{Metadata-Driven Retrieval-Augmented Generation for Financial
Question Answering}, arXiv:2510.24402 (2025). The paper is
catalogued in the program's registry as
\textsc{PAPER\_VERIFIED} with a \emph{preprint under major revision}
caveat. Its algorithmic claims are recorded as the design reference;
no paper-level effectiveness claim is made against this baseline
until the paper is published.

\subsection{Why clean room}
The program's paper-to-spec audit for this baseline records:

\begin{itemize}[leftmargin=1.5em,itemsep=0.15em]
  \item \texttt{repository\_registered}: false
  \item \texttt{repository\_url}: null
  \item \texttt{license\_status}: \textsc{NOT\_APPLICABLE}
  \item \texttt{foreign\_source\_execution\_applicable}: false
\end{itemize}

No foreign source exists for Metadata-RAG. There is no repository to
clone, no code to inspect, no licence to review. The implementation
path is therefore a project-owned paper-derived reimplementation,
structurally similar to the PathRAG comparator but without any
foreign checkout.

\subsection{Fidelity classification}
Under the program's status vocabulary the baseline is classified as:

\begin{center}
\small
\begin{tabular}{ll}
\toprule
\textbf{Field} & \textbf{Value} \\
\midrule
architecture class & \textsc{ARCHITECTURAL\_COMPARATOR} \\
implementation class & \textsc{PAPER\_DERIVED\_REIMPLEMENTATION} \\
fidelity class & \textsc{PAPER\_DERIVED\_COMPARATOR} \\
foreign execution decision & \textsc{NOT\_APPLICABLE} \\
project-owned implementation decision & \textsc{CONDITIONAL\_GO} \\
controlled track eligible & \texttt{true} \\
native track eligible & \texttt{false} \\
native track reason & \textsc{DEFERRED\_WITH\_REASON} \\
\bottomrule
\end{tabular}
\end{center}

The native track is deferred with reason because no foreign source
and no pinned native stack (embedder, tokenizer, storage, metadata
schema) exist. Native-fidelity execution cannot be defined without
inventing components not present in the paper, which the program's
guidance prohibits.

% =====================================================================
\section{Frozen Specification and Its Amendments}\label{sec:frozen}

\subsection{The clean-room specification}
The paper defines a metadata-driven multistage retrieval family but
does not pin its metadata schema, its query-metadata inference rules,
its filtering threshold, its fusion coefficients, or its
evaluation-unit semantics. Every one of those decisions had to be
made explicitly by this project. We made them once, wrote them into
a single frozen specification
(\texttt{metadata\_rag\_clean\_room\_spec.yaml}, SHA-256
\texttt{A02DBA61\ldots}), and treated that file as immutable. Any
subsequent change to a frozen decision is recorded as an explicit
amendment with its own SHA-256 and its own justification.

\subsection{Amendment 01 --- Capitalized-phrase regex symmetry}
\textbf{Motivation.} During pre-code reconnaissance on the synthetic
acceptance fixture, the original capitalized-phrase pattern
\texttt{(?:[A-Z][a-z]+(?:\textbackslash s+[A-Z][a-z]+)+)}
produced the phrase \texttt{"the amber observatory"} from document
text starting with \texttt{"The Amber Observatory records\ldots"},
because the sentence-initial article was absorbed into the matched
span. The corresponding query text produced \texttt{"amber
observatory"} without the article. The document and query metadata
would therefore have failed to align, silently breaking the
architecture-defining symmetry between extraction on the two sides.

\textbf{Patch.} The pattern was replaced by
\texttt{(?<!\textbackslash S)(?!(?:The|A|An)\textbackslash
b)(?:[A-Z][a-z]+(?:\textbackslash s+[A-Z][a-z]+)+)}. A
negative-lookahead guard now excludes sentence-initial articles.

\textbf{Scope.} The amendment changes a regex only. No threshold, no
coefficient, and no outward validation surface changes. The
frozen-spec validator was extended to accept the amended pattern;
the extension is itself a project-owned decision, not a paper claim.

\subsection{Amendment 02 --- In-memory exact dense storage}
\textbf{Motivation.} The clean-room specification originally
declared Qdrant as the experimental storage backend. On review, a
controlled Qdrant deployment would have added an HNSW approximation
risk that would have needed re-freezing for every target corpus
size, an operational server dependency, and an equivalence test
whose outcome on tie-heavy rankings is not guaranteed a priori.
None of this contributes fidelity signal, because the paper pins no
storage backend and the architecture-defining surface is the
metadata layer (schema, extraction, query inference, filtering,
fusion, ranking), not the nearest-neighbour search.

\textbf{Decision.} The experimental backend is
\textsc{IN\_MEMORY\_EXACT\_DENSE}, with
\texttt{replacement\_class} set to
\textsc{SEMANTICALLY\_NEUTRAL\_FOR\_CONTROLLED\_TRACK} and
\texttt{equivalence\_test} set to \textsc{NOT\_REQUIRED}. The
substrate is an exhaustive, unpruned, exact cosine top-$k$ index
implemented as a single NumPy matrix-matrix multiplication. No
approximate search, no quantization, no HNSW, no candidate pruning
is used anywhere in the pipeline.

\textbf{Why the substitution is semantically neutral.} The choice of
exact dense search over an in-memory reference is not a fidelity
compromise; it is the \emph{strongest} semantic guarantee available
in the program. The retrieval signal that reaches the downstream
fusion stage is the exact cosine similarity between the query and
every candidate window. Any approximate backend would only perturb
that signal. By removing the perturbation entirely, the reference
index makes the controlled comparison cleaner, not weaker.

\subsection{Amendment 03 --- Evaluation cutoffs alignment}
\textbf{Motivation.} The program's executed comparison runs
(LongRAG vs.\ My\_RAG) report document-level retrieval metrics at
cutoffs $\{1, 5, 10, 20, 50\}$. The shared
\texttt{common\_retrieval\_evaluation\_contract} additionally
requires cutoff \num{100}. The clean-room spec inherited only
$\{1, 5, 10, 100\}$ from the common contract, which would have made
the Metadata-RAG metrics only partially comparable to the executed
runs.

\textbf{Decision.} The effective cutoff set is the union
$\{1, 5, 10, 20, 50, 100\}$. The primary metric is declared
explicitly as \texttt{macro\_document\_recall\_at\_10}. The
zero-relevance policy is declared explicitly as
\texttt{exclude\_with\_explicit\_count}. No metric is removed, no
threshold is tuned.

\subsection{Amendment 04 --- Oversize document windowing}
\textbf{Motivation.} Real-dataset reconnaissance measured token
counts with the pinned BGE-M3 tokenizer. Every dataset except CLERC
contains canonical documents larger than the BGE-M3 model input of
\num{8192} tokens. Table~\ref{tab:oversize} summarises the observed
distribution.

\begin{table}[h]
\centering
\caption{Token-count distribution per dataset (pinned BGE-M3
tokenizer).}
\label{tab:oversize}
\small
\begin{tabular}{l r r r r r}
\toprule
\textbf{Dataset} & \textbf{Docs} & \textbf{p50} & \textbf{p90}
  & \textbf{max} & \textbf{oversize} \\
\midrule
CLERC              & \num{2756} & \num{383}   & \num{438}
                   & \num{607}   & \num{0.0}\% \\
CUAD               & \num{510}  & \num{7508}  & \num{27913}
                   & \num{77619} & \num{46.5}\% \\
GreekBarBench      & \num{284}  & \num{18307} & \num{30667}
                   & \num{45411} & \num{82.0}\% \\
LaRA               & \num{131}  & \num{62392} & \num{121332}
                   & \num{148481} & \num{100.0}\% \\
\bottomrule
\end{tabular}
\end{table}

Without an explicit oversize policy the encoder would silently
truncate the majority of the source content of three of the four
datasets. Truncation is not an admissible design choice under the
program's controlled protocol because it would make the retrieval
comparison against the executed LongRAG and My\_RAG tracks
meaningless.

\textbf{Decision.} A canonical document whose token count is at or
below \num{8190} content tokens produces exactly one window with
\texttt{window\_index}~$=0$ and \texttt{doc\_id}~$=$
\texttt{lineage\_doc\_id}~$=$ the canonical document identifier. A
canonical document above the budget is split into $N$ contiguous,
exhaustive windows. Each window carries the canonical document's
identity as \texttt{lineage\_doc\_id}, and its own
\texttt{doc\_id} of the form
\texttt{<lineage\_doc\_id>:w<i>}. The metadata record is always
derived from the text of the window itself.

\textbf{Why this is not a fidelity violation.} The paper does not
pin an oversize policy. LongRAG's controlled adaptation uses an
analogous windowing strategy at \num{4096} tokens. The amendment
changes no coefficient, no threshold, no comparison semantics, and
no evaluation unit. It is a no-op on the synthetic acceptance
fixture and on CLERC, which have no oversize documents.

\textbf{Backward compatibility.} The single-window shortcut means
the amendment is byte-identical to the pre-amendment behaviour on
any document that fits within the budget. This was verified at the
test level: the synthetic acceptance run produces identical
\texttt{documents.jsonl} bytes before and after the amendment.

% =====================================================================
\section{Architecture}\label{sec:arch}

\subsection{Overview}
Metadata-RAG is a metadata-driven multistage retrieval architecture.
Its defining surface is a metadata layer that intervenes between the
query and the semantic similarity search. The full offline and online
paths are summarised in Figure~\ref{fig:pipeline}.

\begin{figure}[h]
\centering
\small
\begin{tabular}{c}
\hline
\textbf{Offline}\\
\hline
canonical documents $\to$ deterministic metadata extraction\\
$\to$ windowing (amendment 04) $\to$ indexed windows (state)\\
$\to$ BGE-M3 dense embeddings (Colab, FP16)\\
$\to$ exact in-memory dense index\\
\hline
\textbf{Online}\\
\hline
query text $\to$ metadata extraction (same extractor)\\
$\to$ typed constraints\\
$\to$ candidate windows (match\_ratio $\geq 0.5$)\\
$\to$ exact dense top-$k$ over candidates\\
$\to$ metadata--semantic fusion ($\alpha_m=0.3$, $\alpha_s=0.7$)\\
$\to$ ranked items\\
$\to$ document expansion (first-occurrence dedup by lineage id)\\
\hline
\end{tabular}
\caption{Metadata-RAG pipeline.}
\label{fig:pipeline}
\end{figure}

\subsection{Metadata schema}
The frozen schema has two tiers. The \emph{provenance} tier is
populated from the canonical envelope and is present on every
document window:

\begin{center}
\small
\begin{tabular}{lll}
\toprule
\textbf{Field} & \textbf{Type} & \textbf{Source} \\
\midrule
\texttt{dataset} & categorical & envelope \\
\texttt{doc\_id} & categorical & canonical id \\
\texttt{source\_dataset\_id} & provenance & envelope metadata \\
\texttt{source\_content\_group\_id} & provenance & envelope metadata \\
\texttt{title} & categorical & envelope \\
\bottomrule
\end{tabular}
\end{center}

The \emph{derived} tier is extracted deterministically from the
window text by a small set of frozen regular expressions:

\begin{center}
\small
\begin{tabular}{lll}
\toprule
\textbf{Field} & \textbf{Type} & \textbf{Extraction} \\
\midrule
\texttt{capitalized\_phrases} & entity, multi & regex (amended 01) \\
\texttt{year\_mentions} & temporal, multi & \texttt{\textbackslash b(19|20)\textbackslash d\{2\}\textbackslash b} \\
\texttt{numeric\_tokens} & numeric, multi & \texttt{\textbackslash b\textbackslash d+(\textbackslash .\textbackslash d+)?\textbackslash b} \\
\texttt{char\_length} & numeric, single & \texttt{len(text)} \\
\texttt{word\_count} & numeric, single & whitespace split \\
\texttt{length\_bucket} & categorical, single & word-count thresholds \\
\bottomrule
\end{tabular}
\end{center}

Multi-value fields are deduplicated and sorted. Missing values are
represented as the \texttt{null} token, never as the empty string.

\subsection{The architecture-defining symmetry}
The same extractor is applied to document window text and to query
text. This symmetry is what makes the metadata layer meaningful: the
query side infers the same typed fields that the index side stores.
The query projection reads only \texttt{query\_id} and
\texttt{query\_text}; it never reads \texttt{query.doc\_ids},
\texttt{query.relevant\_ids}, \texttt{query.metadata},
\texttt{query.ground\_truth}, or any other field that could leak
ground truth into retrieval.

\subsection{Candidate filtering}
A typed constraint is emitted for every multi-value derived field
present in the query record. A window's \emph{match ratio} is the
number of constraints it satisfies divided by the total number of
constraints the query emitted. A window is a candidate if its match
ratio is at least the frozen \texttt{min\_match\_ratio} of \num{0.5}.

When the candidate set is empty the pipeline falls back to
\textsc{use\_full\_corpus\_for\_semantic\_retrieval}, sets
\texttt{alpha\_metadata} to \num{0.0}, and flags every emitted item
with \texttt{metadata\_match\_none: true}. The fallback is
deterministic and is reported per query.

\subsection{Semantic retrieval}
Semantic retrieval is dense-only over the pinned BGE-M3 encoder
(\texttt{BAAI/bge-m3}, revision
\texttt{5617a9f6\ldots}), $1024$-dimensional, L2-normalized. Sparse
and ColBERT outputs are disabled. The encoder runs offline against a
local snapshot; network fallback is not permitted. The embedder is
classified \textsc{ARCHITECTURE\_SUPPORTING} with treatment
\textsc{CONTROLLED\_REPLACEMENT}; no native-equivalence claim is
made.

\subsection{Metadata--semantic fusion}
The fused score of a candidate window is
\[
\mathrm{score}(w) \;=\; \alpha_m \cdot \mathrm{match\_ratio}(w)
\;+\; \alpha_s \cdot \max(0,\ \mathrm{cosine}(q, w)),
\]
with $\alpha_m = \num{0.3}$ and $\alpha_s = \num{0.7}$. Both
coefficients are frozen and marked
\texttt{tuning\_prohibited: true}. The paper pins no fusion
formula; this is an explicit project decision and is documented in
the clean-room specification.

Tie-breaking is deterministic: descending fused score, then
descending semantic score, then ascending window \texttt{doc\_id}.

\subsection{Document expansion}
Because one canonical document can be split into multiple windows,
the ranked-item list may contain several windows of the same
document. The expansion stage deduplicates by \texttt{lineage\_doc\_id}
using a first-occurrence rule: the highest-ranked window of each
canonical document determines that document's emitted rank and score.
This matches the LongRAG document-expansion precedent exactly.

% =====================================================================
\section{Component Disposition}\label{sec:disposition}

Every component is classified by its role. The role determines the
admissible treatment.

\begin{center}
\small
\begin{tabular}{lll}
\toprule
\textbf{Component} & \textbf{Role} & \textbf{Treatment} \\
\midrule
metadata schema & defining & preserved \\
metadata extraction & defining & preserved \\
query metadata inference & defining & preserved \\
candidate filtering & defining & preserved \\
semantic retrieval & defining & preserved \\
metadata--semantic fusion & defining & preserved \\
result ranking & defining & preserved \\
source lineage & defining & preserved \\
document expansion & defining & preserved \\
windowing (amend.~04) & defining & preserved \\
embedder & supporting & controlled replacement \\
tokenizer & supporting & matches embedder \\
storage backend & supporting & controlled replacement \\
generation & optional & disabled \\
answer evaluation & optional & disabled \\
llm-as-judge & optional & disabled \\
\bottomrule
\end{tabular}
\end{center}

The design decision to treat storage as supporting rather than
defining is what makes the IN\_MEMORY\_EXACT\_DENSE substitution
admissible under the program's substitution policy. The embedder is
supporting because the paper's contribution is the metadata
multistage retrieval mechanism, not a specialised encoder.

% =====================================================================
\section{Execution Pipeline}\label{sec:pipeline}

Execution is organised into six stages, each implemented as a
resumable, atomic, checksum-closed script.

\begin{center}
\small
\begin{tabular}{c p{4.4cm} p{7.0cm}}
\toprule
\textbf{Stage} & \textbf{Purpose} & \textbf{Outputs} \\
\midrule
L2 & state materialization & \texttt{state-v1/}, \texttt{colab\_bundle/},
     \texttt{colab\_bundle.tar.gz} \\
L3 & Colab BGE-M3 encoding & \texttt{window\_embeddings.jsonl},
     \texttt{query\_embeddings.jsonl}, packaged tar.gz \\
L4 & local import and enrichment & \texttt{embeddings-v1/} with full
     \texttt{MetadataRAGDocumentEmbedding} and
     \texttt{MetadataRAGQueryEmbedding} records \\
L5 & full retrieval & \texttt{retrieval-v1/results.jsonl}, manifest,
     checksums \\
L6 & evaluation & \texttt{evaluation-v1/query\_metrics.jsonl},
     \texttt{aggregate\_metrics.json} \\
L7 & cross-track comparison & \texttt{baseline\_audit/reports/} JSON
     and Markdown \\
\bottomrule
\end{tabular}
\end{center}

Every stage writes a manifest with the SHA-256 of every artifact it
produces, and every later stage verifies the hashes of the artifacts
it consumes. A silent change to any upstream artifact produces a
hard failure at the next stage.

\subsection{Determinism and reproducibility}
All stages are byte-deterministic given their inputs. Two
consecutive runs of the synthetic acceptance produce identical
artifacts at every stage. Two consecutive runs of the full retrieval
on the same embeddings produce identical ranked lists. The only
source of nondeterminism in the full pipeline is FP16 rounding in
the Colab encoder, which is bounded and is not claimed to be
bit-reproducible across GPUs.

\subsection{Leakage prevention}
The pipeline's query projection, the retrieval stage, and the
retrieval output all read only the fields they are permitted to
read. The evaluator is an independent process that reads retrieval
results and canonical ground truth, and it is the only stage that
ever sees ground truth. The Colab encoder never receives metadata
records, never receives ground truth, and never sees query text
beyond the query window text it is asked to embed.

% =====================================================================
\section{Real-Dataset Execution}\label{sec:real}

\subsection{Datasets and scope}
The executed scope is the same four datasets as the LongRAG and
My\_RAG tracks: CLERC, CUAD, GreekBarBench, and LaRA.
Table~\ref{tab:datasets} reports the canonical scope and the
materialized window count for Metadata-RAG.

\begin{table}[h]
\centering
\caption{Executed scope and materialized windows.}
\label{tab:datasets}
\small
\begin{tabular}{l r r r r}
\toprule
\textbf{Dataset} & \textbf{Documents} & \textbf{Queries}
  & \textbf{Windows} & \textbf{Expansion factor} \\
\midrule
CLERC              & \num{2756} & \num{1000}  & \num{2756}
                   & \num{1.00} \\
CUAD               & \num{510}  & \num{20910} & \num{1017}
                   & \num{1.99} \\
GreekBarBench      & \num{284}  & \num{284}   & \num{792}
                   & \num{2.79} \\
LaRA               & \num{131}  & \num{2326}  & \num{1160}
                   & \num{8.85} \\
\midrule
\textbf{Total}     & \num{3681} & \num{24520} & \num{5725}
                   & --- \\
\bottomrule
\end{tabular}
\end{table}

The window counts match the pre-execution token-budget projection
exactly, which is itself evidence that the amendment-04 windowing
policy was implemented correctly.

\subsection{Embedding}
Every window and every query was embedded with the pinned BGE-M3
encoder in FP16 on a Google Colab Tesla T4. The Colab-side script
records only the \texttt{(doc\_id, dense\_vector)} and
\texttt{(query\_id, dense\_vector)} pairs. The local import stage
reattaches metadata from the frozen state and enforces a strict
bijection: every state window must have exactly one embedding, every
embedding must refer to a state window, and every query embedding
must correspond to a query text in the bundle.

\subsection{Retrieval configuration}
Retrieval used the frozen configuration: \texttt{alpha\_metadata}
$=0.3$, \texttt{alpha\_semantic} $=0.7$,
\texttt{min\_match\_ratio} $=0.5$, \texttt{document\_limit} $=100$.
The configuration object refuses construction if any of these values
is changed.

\subsection{Observed timings}
\begin{center}
\small
\begin{tabular}{l r r}
\toprule
\textbf{Dataset} & \textbf{Queries} & \textbf{Retrieval wall time} \\
\midrule
CLERC              & \num{1000}  & \SI{169.5}{\second} \\
CUAD               & \num{20910} & \SI{255.9}{\second} \\
GreekBarBench      & \num{284}   & \SI{3.2}{\second} \\
LaRA               & \num{2326}  & \SI{29.9}{\second} \\
\bottomrule
\end{tabular}
\end{center}

The corpus-size dependence is dominated by the candidate-set size:
CUAD has \num{20910} queries each scoped to its own candidate set,
and CLERC has the largest full-corpus fallback, so both dominate the
wall time even though their window counts are moderate.

% =====================================================================
\section{Results}\label{sec:results}

\subsection{Primary metric}
The primary metric of the program is macro-averaged document-level
\texttt{Recall@10}. Table~\ref{tab:primary} reports it per dataset
and as a macro across the four datasets.

\begin{table}[h]
\centering
\caption{Metadata-RAG document \texttt{Recall@10}.}
\label{tab:primary}
\small
\begin{tabular}{l r r r}
\toprule
\textbf{Dataset} & \textbf{Queries evaluated}
  & \textbf{Queries excluded (0 relevant)}
  & \textbf{Recall@10} \\
\midrule
CLERC              & \num{1000} & \num{0}     & \num{0.2726} \\
CUAD               & \num{6702} & \num{14208} & \num{0.0294} \\
GreekBarBench      & \num{284}  & \num{0}     & \num{0.2711} \\
LaRA               & \num{2326} & \num{0}     & \num{0.8874} \\
\midrule
\textbf{Macro}     & ---        & ---         & \num{0.3651} \\
\bottomrule
\end{tabular}
\end{table}

\subsection{Full metric set}
Table~\ref{tab:full} reports the complete macro-averaged metric set
at every cutoff.

\begin{table}[h]
\centering
\caption{Metadata-RAG macro-averaged metrics across four datasets.}
\label{tab:full}
\small
\begin{tabular}{l r | l r}
\toprule
\textbf{Metric} & \textbf{Value}
  & \textbf{Metric} & \textbf{Value} \\
\midrule
Hit@1     & \num{0.2313} & Recall@1   & \num{0.2026} \\
Hit@5     & \num{0.3389} & Recall@5   & \num{0.2999} \\
Hit@10    & \num{0.4085} & Recall@10  & \num{0.3651} \\
Hit@20    & \num{0.4789} & Recall@20  & \num{0.4316} \\
Hit@50    & \num{0.5920} & Recall@50  & \num{0.5478} \\
Hit@100   & \num{0.6570} & Recall@100 & \num{0.6121} \\
Precision@10 & \num{0.0482} & MRR@10  & \num{0.2797} \\
nDCG@10   & \num{0.2883} & MRR@50     & \num{0.2882} \\
nDCG@50   & \num{0.3319} & & \\
\bottomrule
\end{tabular}
\end{table}

\subsection{Per-dataset metric table}
Table~\ref{tab:perds} reports the per-dataset Hit@1, Hit@10, Recall@1,
Recall@10, Recall@100, MRR@10, and nDCG@10.

\begin{table}[h]
\centering
\caption{Metadata-RAG per-dataset primary metrics.}
\label{tab:perds}
\small
\begin{tabular}{l r r r r r r r}
\toprule
\textbf{Dataset} & \textbf{H@1} & \textbf{H@10} & \textbf{R@1}
  & \textbf{R@10} & \textbf{R@100} & \textbf{MRR@10} & \textbf{nDCG@10} \\
\midrule
CLERC              & \num{0.1870} & \num{0.4460}
                   & \num{0.0723} & \num{0.2726}
                   & \num{0.5322} & \num{0.2626} & \num{0.2199} \\
CUAD               & \num{0.0034} & \num{0.0279}
                   & \num{0.0034} & \num{0.0294}
                   & \num{0.2125} & \num{0.0090} & \num{0.0137} \\
GreekBarBench      & \num{0.0176} & \num{0.2711}
                   & \num{0.0176} & \num{0.2711}
                   & \num{0.7183} & \num{0.0756} & \num{0.1204} \\
LaRA               & \num{0.7171} & \num{0.8874}
                   & \num{0.7171} & \num{0.8874}
                   & \num{0.9854} & \num{0.7715} & \num{0.7993} \\
\bottomrule
\end{tabular}
\end{table}

% =====================================================================
\section{Architectural Observation: The Metadata Stage Did Not
Activate on Three of Four Corpora}\label{sec:observation}

\subsection{The measurement}
The single most important observation of this track is not the
aggregate macro metric. It is the per-query fallback rate of the
metadata-filtering stage. Table~\ref{tab:fallback} reports it.

\begin{table}[h]
\centering
\caption{Metadata-stage fallback rate per dataset.}
\label{tab:fallback}
\small
\begin{tabular}{l r r r}
\toprule
\textbf{Dataset} & \textbf{Queries} & \textbf{Fell back}
  & \textbf{Fallback rate} \\
\midrule
CLERC              & \num{1000}  & \num{999}   & \num{99.9}\% \\
CUAD               & \num{20910} & \num{20400} & \num{97.6}\% \\
GreekBarBench      & \num{284}   & \num{262}   & \num{92.3}\% \\
LaRA               & \num{2326}  & \num{784}   & \num{33.7}\% \\
\bottomrule
\end{tabular}
\end{table}

On three of the four datasets the metadata filter is inactive for
essentially the entire query set. Under the frozen retrieval
semantics this means that the architecture executed as a pure
dense-only BGE-M3 retriever on those three datasets, and as a
metadata-filtered dense retriever only on LaRA. The measured
retrieval quality is therefore the quality of a dense-only retriever
on three of the four datasets.

\subsection{The mechanism}
The mechanism is the interaction of two frozen decisions with the
properties of the corpora:

\begin{enumerate}[leftmargin=1.5em,itemsep=0.15em]
  \item The frozen \texttt{min\_match\_ratio} of \num{0.5} requires
    a candidate window to satisfy at least half of the query's
    emitted constraints.
  \item The frozen extractor emits one typed constraint per detected
    multi-value field value in the query text.
\end{enumerate}

When a query produces many constraints, the second requirement
becomes unattainable. We observed exactly this on CLERC.

\subsection{Why CLERC fails}
CLERC query texts are long judicial opinions, averaging hundreds of
tokens and containing dozens of multi-word capitalized phrases. The
extractor emitted between \num{65} and \num{137} constraints per
query, with a median well above the $60$ threshold observed in the
diagnostic histogram. Under a $\num{0.5}$ ratio, a window would have
to satisfy \num{33}--\num{69} distinct constraints simultaneously,
which no canonical document achieves. The candidate set is empty on
\num{99.9}\% of CLERC queries and the fallback triggers.

\subsection{Why CUAD fails}
CUAD query texts are templated prompts of the form
\emph{``Highlight the parts (if any) of this contract related to
`X' that should be reviewed by a lawyer.''} The extractor emits
\num{0} or \num{1} constraints depending on whether the templated
slot contains a capitalized phrase. When \num{0} constraints are
emitted, the ratio is undefined and the fallback triggers by
construction. \num{65.9}\% of CUAD queries produced zero
constraints; \num{34.1}\% produced exactly one. In both cases, the
frozen ratio of \num{0.5} cannot be satisfied.

\subsection{Why GreekBarBench fails}
GreekBarBench query texts are in Greek. The frozen extractor uses
ASCII ranges only
(\texttt{[A-Z][a-z]}, \texttt{\textbackslash b(19|20)\textbackslash
d\{2\}\textbackslash b}, etc.) and was deliberately not extended to
Unicode capitals in any amendment. On Greek text the extractor emits
almost no constraints: \num{89.8}\% of queries produced zero
constraints, and the fallback triggers.

This is a bounded and documented limitation of the frozen spec,
not a runtime bug. The spec's schema version and the extractor
version are both pinned, and the limitation is stated in the
extractor module's own docstring.

\subsection{Why LaRA works}
LaRA query texts are English natural questions of the form
\emph{``How does Lucy's decision to end her engagement with Cecil
differ from her response to Mr.\ Emerson's encouragement about her
feelings for George?''} The extractor emits \num{1} to \num{5}
constraints per query. A $\num{0.5}$ ratio then requires a window to
satisfy \num{1} to \num{3} constraints, which is achievable. The
metadata filter activates on \num{66.3}\% of LaRA queries.

\subsection{Scientific interpretation}
The observation can be stated as a design-level finding:

\begin{quote}
\emph{A metadata-driven retrieval architecture with a frozen
minimum-match ratio produces a filter that is only active on corpora
whose query texts emit a constrained number of typed metadata
signals. When query texts are long, templated, or written in a
language the frozen extractor cannot tokenise, the filter
degenerates to a dense-only pipeline and the metadata stage
contributes nothing.}
\end{quote}

This is not a defect of the paper's algorithmic family. It is a
property of the interaction between the family and any fixed
configuration of its parameters. In a production deployment the
parameter would be tuned per corpus; in a controlled comparison it
must be frozen, and the frozen value produces the observed
behaviour. The comparison therefore measures the architecture's
behaviour under a fixed, pre-registered configuration, which is
exactly what a controlled architectural comparison is meant to do.

\subsection{No post-hoc tuning}
The program's master guidance prohibits post-hoc threshold tuning.
The frozen \texttt{min\_match\_ratio} of \num{0.5} was declared
before execution, marked \texttt{tuning\_prohibited: true} in the
clean-room specification, and remained unchanged throughout. The
observation is recorded in the readiness declaration's
\texttt{architectural\_observation} block, and the flag
\texttt{post\_hoc\_tuning\_applied} is set to \texttt{false}.

% =====================================================================
\section{Three-Track Comparison}\label{sec:comparison}

\subsection{Setup}
The comparison covers the three tracks that have been executed on
the identical four datasets under the identical shared protocol:
LongRAG (controlled adaptation), Metadata-RAG (paper-derived
comparator), and My\_RAG (project-owned reference). The comparison
reads only per-dataset aggregate metrics of each track; it never
touches per-query data.

\subsection{Primary metric comparison}
Table~\ref{tab:three} reports the per-dataset and macro primary
metric.

\begin{table}[h]
\centering
\caption{Macro-averaged document \texttt{Recall@10} across tracks.}
\label{tab:three}
\small
\begin{tabular}{l r r r}
\toprule
\textbf{Dataset} & \textbf{LongRAG} & \textbf{Metadata-RAG}
  & \textbf{My\_RAG} \\
\midrule
CLERC              & \num{0.0786} & \num{0.2726} & \num{0.3546} \\
CUAD               & \num{0.0324} & \num{0.0294} & \num{0.0431} \\
GreekBarBench      & \num{0.2606} & \num{0.2711} & \num{0.2852} \\
LaRA               & \num{0.8861} & \num{0.8874} & \num{0.9368} \\
\midrule
\textbf{Macro}     & \num{0.3144} & \num{0.3651} & \num{0.4049} \\
\bottomrule
\end{tabular}
\end{table}

Metadata-RAG is \num{16.1}\% above LongRAG and \num{9.8}\% below
My\_RAG on the primary metric.

\subsection{Full macro metric table}
Table~\ref{tab:threefull} reports the complete macro metric set.

\begin{table}[h]
\centering
\caption{Macro metric set across tracks.}
\label{tab:threefull}
\small
\begin{tabular}{l r r r}
\toprule
\textbf{Metric} & \textbf{LongRAG} & \textbf{Metadata-RAG}
  & \textbf{My\_RAG} \\
\midrule
Hit@1     & \num{0.1957} & \num{0.2313} & \num{0.2553} \\
Hit@5     & \num{0.2606} & \num{0.3389} & \num{0.3846} \\
Hit@10    & \num{0.3183} & \num{0.4085} & \num{0.4520} \\
Hit@20    & \num{0.3881} & \num{0.4789} & \num{0.5416} \\
Hit@50    & \num{0.4929} & \num{0.5920} & \num{0.6658} \\
Recall@1  & \num{0.1905} & \num{0.2026} & \num{0.2180} \\
Recall@5  & \num{0.2569} & \num{0.2999} & \num{0.3394} \\
Recall@10 & \num{0.3144} & \num{0.3651} & \num{0.4049} \\
Recall@20 & \num{0.3833} & \num{0.4316} & \num{0.4938} \\
Recall@50 & \num{0.4870} & \num{0.5478} & \num{0.6210} \\
Precision@10 & \num{0.0362} & \num{0.0482} & \num{0.0557} \\
MRR@10    & \num{0.2268} & \num{0.2797} & \num{0.3123} \\
MRR@50    & \num{0.2349} & \num{0.2882} & \num{0.3225} \\
nDCG@10   & \num{0.2476} & \num{0.2883} & \num{0.3214} \\
nDCG@50   & \num{0.2880} & \num{0.3319} & \num{0.3733} \\
\bottomrule
\end{tabular}
\end{table}

\subsection{Interpretation of the ranking}
Metadata-RAG sits between LongRAG and My\_RAG on the primary metric.
Two architectural facts explain this ordering.

\paragraph{Why Metadata-RAG beats LongRAG.}
LongRAG groups adjacent documents into long retrieval units of
$\approx\num{4096}$ tokens. On CLERC, whose documents average a few
hundred tokens, this packs roughly eleven documents per unit and
dilutes a single relevant document by ten irrelevant ones. Metadata-RAG
does not group; its retrieval unit is a single document window. Even
with the metadata stage inactive, the finer unit preserves document
identity in the ranking. This is what produces the \num{3.47}$\times$
relative advantage on CLERC.

\paragraph{Why My\_RAG beats Metadata-RAG.}
My\_RAG combines three architectural mechanisms that Metadata-RAG
does not have: hierarchical chunking that never crosses
semantic-parent boundaries, a learned-sparse lexical signal fused
with the dense signal by Reciprocal Rank Fusion, and a
parent-aggregation rule that rewards parents with multiple strong
child chunks. On a corpus where the metadata filter is inactive,
Metadata-RAG is a dense-only retriever against a hybrid retriever
with those additional mechanisms. The gap of \num{9.8}\% on the
primary metric is consistent with that difference.

\paragraph{Why the LaRA gap is small.}
On LaRA, the metadata filter activates on \num{66.3}\% of queries
and Metadata-RAG's lead over LongRAG narrows to \num{0.15}\%
absolute. This is consistent with the LaRA corpus being a
near-document-level retrieval problem for both architectures: when
the retrieval unit is already close to a full document, unit
granularity has little room to differentiate the two.

% =====================================================================
\section{Threats to Validity}\label{sec:threats}

\subsection{Paper-status caveat}
The source paper is a preprint under major revision. No paper-level
effectiveness claim is made against this baseline. The paper's
algorithmic claims are the design reference only.

\subsection{No native-fidelity track}
The architecture has no native-fidelity execution. No foreign source
exists and the paper pins no native embedder, tokenizer, storage
backend, or metadata schema. The native track is
\textsc{DEFERRED\_WITH\_REASON}.

\subsection{Controlled adaptations}
Three controlled adaptations are present and are documented:
\begin{itemize}[leftmargin=1.5em,itemsep=0.15em]
  \item the embedder is BGE-M3 dense-only, FP16, pinned revision;
  \item the storage backend is IN\_MEMORY\_EXACT\_DENSE;
  \item an oversize-document windowing policy (amendment 04) is
    applied to documents that exceed the BGE-M3 input budget.
\end{itemize}
None of these changes the architecture-defining retrieval logic.

\subsection{Frozen threshold, not a tuned threshold}
The \texttt{min\_match\_ratio} of \num{0.5} was pre-registered and
never tuned. The fallback rate that results from this choice is a
property of the frozen design and is reported as-is. A production
deployment would tune the threshold per corpus; the controlled
comparison is not permitted to.

\subsection{Single-run execution}
Both this track and its comparison tracks are executed once,
deterministically, on the full query set. No bootstrap confidence
intervals and no paired significance tests are reported. The
observed differences are descriptive.

\subsection{Precision}
FP16 embeddings are used. FP32 equivalence is not claimed.

\subsection{Ground-truth projection}
All three tracks are evaluated at the canonical-document level.
Datasets such as CUAD and GreekBarBench have span-level or
clause-level annotations that could arguably be evaluated at finer
granularity; the program's shared protocol commits to document level
and the reported numbers reflect that choice.

% =====================================================================
\section{Reproducibility}\label{sec:repro}

\subsection{Commit chain}
The Metadata-RAG track closure corresponds to the following commits,
in order:

\begin{center}
\small
\begin{tabular}{l l}
\toprule
\textbf{Commit} & \textbf{Stage} \\
\midrule
\texttt{c298429} & clean-room spec + controlled-track pipeline + synthetic readiness \\
\texttt{a91bed7} & oversize document windowing (amendment 04) \\
\texttt{bf63ee2} & state materialization + Colab bundle prep \\
\texttt{cabf1e8} & L2--L6 scripts (Colab encode, import, retrieval, evaluation, comparison) \\
\texttt{8e23fe4} & evaluator dedup fix + legacy-format comparison \\
\texttt{d7186fe} & readiness declaration REAL\_DATASET\_EXECUTED \\
\bottomrule
\end{tabular}
\end{center}

\subsection{Artifact inventory}
\begin{center}
\small
\begin{tabular}{l l}
\toprule
\textbf{Artifact} & \textbf{Location} \\
\midrule
Frozen spec & \texttt{baseline\_audit/metadata\_rag\_clean\_room\_spec.yaml} \\
Amendments 01--04 & \texttt{baseline\_audit/metadata\_rag\_clean\_room\_spec\_amendment\_*.yaml} \\
Readiness declaration & \texttt{baseline\_audit/metadata\_rag\_readiness\_declaration.yaml} \\
State, embeddings, retrieval, evaluation & \texttt{outputs/external\_metadata\_rag/<ds>/} \\
Comparison report & \texttt{baseline\_audit/reports/metadata\_rag\_vs\_baselines.\{json,md\}} \\
\bottomrule
\end{tabular}
\end{center}

\subsection{Determinism}
All artifacts are byte-deterministic given their inputs. Every stage
writes a manifest with SHA-256 hashes, and later stages verify the
hashes of the artifacts they consume. Silent drift is detected as a
hard failure.

\subsection{Test coverage}
\num{342} tests across nineteen test modules. All pass under ruff,
mypy strict, and pytest. The suite covers the metadata extractor,
the adapter, the persistence layer, the embedder contract, the
reference index, the retrieval pipeline, the expansion stage, the
evaluation-case builder, the windowing module, the synthetic
acceptance runner, the Colab encoder, the import-enrich stage, the
full-retrieval driver, the evaluator, the comparison script, and the
readiness declaration validator.

% =====================================================================
\section{Conclusion}\label{sec:conclusion}

We have described, frozen, implemented, executed, and evaluated
Metadata-RAG, a paper-derived reimplementation of a
metadata-driven multistage retrieval architecture, on four
heterogeneous canonical corpora. The architecture is complete,
deterministic, fully validated, and passes \num{342} tests. Its
macro-averaged document \texttt{Recall@10} across the four datasets
is \num{0.3651}, between LongRAG's \num{0.3144} and My\_RAG's
\num{0.4049}.

The substantive finding of this track is not the aggregate metric but
its decomposition. Under the frozen \texttt{min\_match\_ratio} of
\num{0.5}, the metadata-filtering stage failed to activate on three
of the four datasets, and the architecture executed as a pure
dense-only retriever on those three. We documented this behaviour
with full transparency, we identified its mechanism, and we made no
post-hoc adjustment of the frozen specification. The observation is
architecturally useful: it demonstrates that metadata-driven
retrieval requires the query text to produce a constrained number of
typed signals, and that a frozen filter threshold can be inactive on
corpora whose query texts are long, templated, or written in a
language the extractor does not cover.

The track is declared closed at commit \texttt{d7186fe} with status
\textsc{REAL\_DATASET\_EXECUTED}. The next action of the program is
the comparative report writing that will place the three tracks
side by side and state the architectural observations that emerge.

% =====================================================================
\appendix

\section{Appendix A: The Clean-Room Spec at a Glance}
\begin{lstlisting}[language=yaml]
architecture:
  architecture_id: metadata_rag
  fidelity_class: PAPER_DERIVED_COMPARATOR
  implementation_class: PAPER_DERIVED_REIMPLEMENTATION
tracks:
  controlled: { eligible: true }
  native:     { eligible: false, reason: DEFERRED_WITH_REASON }
metadata_schema:
  schema_version: "0.1.0"
  schema_status: RESOLVED
query_metadata_inference:
  classification: DETERMINISTIC_PROJECT_OWNED
  llm_required: false
  query_projection:
    allowed_fields: [query_id, query_text]
    prohibited_fields:
      - doc_ids
      - relevant_ids
      - answers
      - choice
      - target_content_group_id
      - ground_truth
      - spans
      - metadata
candidate_filtering:
  min_match_ratio: 0.5
  min_match_ratio_frozen: true
  min_match_ratio_tuning_prohibited: true
metadata_semantic_fusion:
  coefficients:
    alpha_metadata: 0.3
    alpha_semantic: 0.7
    frozen: true
    tuning_prohibited: true
semantic_retrieval:
  signal: dense_only
  embedder:
    model_id: BAAI/bge-m3
    model_revision: 5617a9f61b028005a4858fdac845db406aefb181
    dense_dimension: 1024
    dense_normalized: true
    network_fallback_allowed: false
storage:
  experimental_backend: IN_MEMORY_EXACT_DENSE
  replacement_class: SEMANTICALLY_NEUTRAL_FOR_CONTROLLED_TRACK
  equivalence_test: NOT_REQUIRED
  approximate_search_allowed: false
  quantization_allowed: false
dataset_compatibility:
  cutoffs: [1, 5, 10, 20, 50, 100]
  primary_metric: macro_document_recall_at_10
  zero_relevance_policy: exclude_with_explicit_count
claims:
  faithful_reproduction: false
  paper_effectiveness_reproduced: false
  native_equivalence: false
  controlled_adaptation: true
  paper_derived_comparator: true
  real_dataset_validated: true
  publishable_result: false
  synthetic_validation_is_paper_level_evidence: false
\end{lstlisting}

\section{Appendix B: The Four Amendments at a Glance}
\begin{center}
\small
\begin{tabular}{p{0.6cm} p{3.6cm} p{5.0cm} p{4.5cm}}
\toprule
\textbf{\#} & \textbf{Kind} & \textbf{Change} & \textbf{Reason} \\
\midrule
01 & regex correctness &
  capitalized-phrase pattern now excludes sentence-initial
  articles \texttt{The|A|An} &
  document/query extraction symmetry would otherwise be broken
  silently on the synthetic fixture \\
02 & storage backend &
  Qdrant replaced by IN\_MEMORY\_EXACT\_DENSE &
  paper pins no backend; exact dense index removes approximation
  risk without fidelity cost \\
03 & evaluation profile &
  cutoffs extended to $\{1,5,10,20,50,100\}$, primary metric
  declared explicitly &
  align with the executed LongRAG and My\_RAG runs \\
04 & oversize handling &
  fixed-size windowing at \num{8190} content tokens using the
  pinned BGE-M3 tokenizer &
  three of the four datasets contain documents that would otherwise
  be truncated by one to two orders of magnitude \\
\bottomrule
\end{tabular}
\end{center}

\section{Appendix C: The Fallback Rates in Detail}
\begin{center}
\small
\begin{tabular}{l r r r r}
\toprule
\textbf{Dataset} & \textbf{Queries} & \textbf{Fallback}
  & \textbf{Mean candidate count} & \textbf{Max candidate count} \\
\midrule
CLERC              & \num{1000}  & \num{99.9}\%
                   & \num{2753.2} & \num{2756} \\
CUAD               & \num{20910} & \num{97.6}\%
                   & \num{992.3}  & \num{1017} \\
GreekBarBench      & \num{284}   & \num{92.3}\%
                   & \num{749.0}  & \num{792} \\
LaRA               & \num{2326}  & \num{33.7}\%
                   & \num{693.8}  & \num{1160} \\
\bottomrule
\end{tabular}
\end{center}

The mean candidate count is close to the full corpus size on every
dataset, which confirms that when the fallback triggers it triggers
\emph{to the full corpus}, not to an empty set. This is the frozen
fallback behaviour declared in the clean-room specification.

\end{document}